\chapter{Radon--Nikodym の定理}
\label{ch:radon-nikodym}

Radon--Nikodym の定理は，一つの測度を別の測度に関する密度で表す定理である．
第\ref{ch:disintegration}章では，この密度を使って各 $y$ における積分の値を定める．
主張は北別府悠『最適輸送理論入門〜兼非標準的関数解析入門〜』
Corollary 4.6.2（p.~57）に従い，証明には正集合の補題を用いる．

以下，$(\Omega,\mathcal{A})$ を可測空間とする．

\section{絶対連続性と定理の主張}

\begin{definition}{測度の絶対連続性（北別府 p.~57）}{meas-absolute-continuity}
  $(\Omega,\mathcal{A})$ 上の有限測度 $\nu,\alpha$ を考える．
  $\alpha$ が $\nu$ に\textbf{絶対連続}であるとは，
  任意の $\varepsilon>0$ に対してある $\delta>0$ が存在し，
  \[
    \nu(B)<\delta\quad\Longrightarrow\quad\alpha(B)<\varepsilon
    \qquad(B\in\mathcal{A})
  \]
  が成り立つことをいう．このとき $\alpha\ll\nu$ と書く．
\end{definition}

\begin{remark}{零集合による特徴付けとの関係}{rn-null-characterization}
  文献によっては（例えば Folland (1999) §3.2），
  「$\nu(B)=0$ ならば $\alpha(B)=0$」を絶対連続性の定義とする．
  有限測度の場合，この二つは同値である（定理~\ref{thm:meas-radon-nikodym}）．
\end{remark}

\begin{theorem}{Radon--Nikodym の定理（北別府 Cor 4.6.2）}{meas-radon-nikodym}
  $(\Omega,\mathcal{A})$ 上の有限測度 $\nu,\alpha$ に対し，次の三つは同値である．
  \begin{enumerate}
    \item $\alpha\ll\nu$．
    \item すべての $B\in\mathcal{A}$ に対して，$\nu(B)=0$ ならば $\alpha(B)=0$ である．
    \item 非負の $\nu$-可積分な可測関数 $h\colon\Omega\to[0,\infty)$ が存在し，
      \[
        \alpha(B)=\int_Bh\,\d\nu\qquad(B\in\mathcal{A})
      \]
      が成り立つ．
  \end{enumerate}
  条件 3 の $h$ は $\nu$-a.e. の意味で一意であり，
  \textbf{Radon--Nikodym 微分}とよんで $\d\alpha/\d\nu$ と書く．
\end{theorem}

証明の中心は $2\Rightarrow3$ である．
$\alpha$ を超えない積分を与える密度の中から，積分値が最大のものを取り，
残りの測度が $0$ であることを示す．その準備として二つの補題を証明する．

\section{証明の準備}

\begin{lemma}{密度の一意性}{rn-density-uniqueness}
  $\nu$ を $(\Omega,\mathcal{A})$ 上の有限測度とし，
  $h,k\colon\Omega\to\R$ を $\nu$-可積分な可測関数とする．すべての $B\in\mathcal{A}$ で
  \[
    \int_Bh\,\d\nu=\int_Bk\,\d\nu
  \]
  が成り立つならば，$h=k$ が $\nu$-a.e. に成り立つ．
\end{lemma}

\begin{proof}
  $B:=\{h>k\}$ では $h-k>0$ である一方，$\int_B(h-k)\,\d\nu=0$ なので $\nu(B)=0$ である．
  $h,k$ を入れ替えて $\nu(\{k>h\})=0$ も得る．
\end{proof}

次の補題は，測度の差が正になる部分を取り出すために用いる．
証明は Hahn の補題と同じく，値が負になる部分集合を順に取り除く構成による
（Royden--Fitzpatrick \emph{Real Analysis}，第4版）．

\begin{lemma}{正集合の存在}{rn-positive-set}
  $\alpha,\beta$ を $(\Omega,\mathcal{A})$ 上の有限測度とし，$\lambda:=\alpha-\beta$ とおく．
  $\lambda(\Omega)>0$ ならば，ある $P\in\mathcal{A}$ で
  \[
    \lambda(P)>0,\qquad
    \lambda(C)\geq0\quad(C\in\mathcal{A},\ C\subset P)
  \]
  をみたすものが存在する．
\end{lemma}

\begin{proof}
  互いに素な可測集合 $C_k$ に対し
  \[
    \sum_k\abs{\lambda(C_k)}\leq\alpha(\Omega)+\beta(\Omega)<\infty
  \]
  なので，$\lambda(\bigcup_kC_k)=\sum_k\lambda(C_k)$ が成り立つ．

  $A_0:=\Omega$ とし，$A_{k-1}$ が定まったとき
  \[
    m_k:=\inf\{\lambda(C)\mid C\in\mathcal{A},\ C\subset A_{k-1}\}
  \]
  とおく．$-\beta(\Omega)\leq m_k\leq0$ である．
  $m_k=0$ ならば $P:=A_{k-1}$ として終了する．
  $m_k<0$ ならば，可測集合 $C_k\subset A_{k-1}$ を
  $\lambda(C_k)<m_k/2$ となるように取り，$A_k:=A_{k-1}\setminus C_k$ とする．
  負の値をもつ集合だけを取り除くので，$\lambda(A_k)\geq\lambda(\Omega)>0$ であり，
  操作が止まれば求める $P$ が得られる．

  止まらない場合は $P:=\Omega\setminus\bigcup_kC_k$ とおく．
  $C_k$ は互いに素なので
  \[
    \lambda(P)=\lambda(\Omega)-\sum_k\lambda(C_k)\geq\lambda(\Omega)>0.
  \]
  また，絶対収束より $\lambda(C_k)\to0$ であり，$2\lambda(C_k)<m_k\leq0$ から $m_k\to0$ である．
  可測集合 $C\subset P$ はすべての $A_{k-1}$ に含まれるので，
  $\lambda(C)\geq m_k$ がすべての $k$ で成り立ち，$\lambda(C)\geq0$ を得る．
\end{proof}

\section{定理の証明と積分への応用}

\begin{proof}[定理~\ref{thm:meas-radon-nikodym}の証明]
  \textbf{$1\Rightarrow2$．}
  $\nu(B)=0$ ならば，絶対連続性より任意の $\varepsilon>0$ に対して
  $\alpha(B)<\varepsilon$ となるので，$\alpha(B)=0$ である．

  \textbf{$3\Rightarrow1$．}
  $\varepsilon>0$ を取る．$h$ の可積分性と単調収束定理（定理~\ref{thm:meas-mct}）より，
  $\int_{\{h>n\}}h\,\d\nu\to0$ である．
  この積分が $\varepsilon/2$ 未満となる $n\in\N$ を取り，$\delta:=\varepsilon/(2n)$ とすれば，
  $\nu(B)<\delta$ のとき
  \[
    \alpha(B)=\int_Bh\,\d\nu
    \leq\int_{\{h>n\}}h\,\d\nu+n\,\nu(B)<\varepsilon.
  \]

  \textbf{$2\Rightarrow3$．}
  \[
    \mathcal{G}:=\left\{g\colon\Omega\to[0,\infty]\ \middle|\
    g\text{ は可測で，すべての }B\in\mathcal{A}\text{ について }\int_Bg\,\d\nu\leq\alpha(B)\right\}
  \]
  とおく．$0\in\mathcal{G}$ である．

  \emph{(a) 積分値が最大の密度を作る．}
  $g_1,g_2\in\mathcal{G}$ とし，$B\in\mathcal{A}$ を
  $B_1:=B\cap\{g_1\geq g_2\}$，$B_2:=B\setminus B_1$ に分けると
  \[
    \int_B\max\{g_1,g_2\}\,\d\nu
    =\int_{B_1}g_1\,\d\nu+\int_{B_2}g_2\,\d\nu
    \leq\alpha(B_1)+\alpha(B_2)=\alpha(B).
  \]

  $s:=\sup_{g\in\mathcal{G}}\int_\Omega g\,\d\nu\leq\alpha(\Omega)<\infty$ とおく．
  $\int_\Omega g_n\,\d\nu\to s$ となる $g_n\in\mathcal{G}$ を取り，
  $\tilde g_n:=\max\{g_1,\dots,g_n\}$ とおくと，前段より $\tilde g_n\in\mathcal{G}$ であり，
  $\tilde g_n$ は単調増加である．その各点極限を $g$ とすると，
  単調収束定理（定理~\ref{thm:meas-mct}）より，すべての $B$ で
  $\int_Bg\,\d\nu=\lim_n\int_B\tilde g_n\,\d\nu\leq\alpha(B)$ となり，
  $g\in\mathcal{G}$ かつ $\int_\Omega g\,\d\nu=s$ である．
  $\int_\Omega g\,\d\nu<\infty$ より $\nu(\{g=\infty\})=0$ なので，
  この零集合上で $g$ の値を $0$ に置き換えた関数を $h$ とすれば，
  $h$ は実数値で，$h\in\mathcal{G}$，$\int_\Omega h\,\d\nu=s$ である．

  \emph{(b) 残りの測度は $0$ である．}
  \[
    \beta(B):=\alpha(B)-\int_Bh\,\d\nu\qquad(B\in\mathcal{A})
  \]
  とおくと，$h\in\mathcal{G}$ より $\beta\geq0$ であり，$\beta$ は有限測度である
  （可算加法性は単調収束定理による）．$\beta(\Omega)=0$ を示せばよい．

  $\beta(\Omega)>0$ と仮定する．
  $t>0$ を $\beta(\Omega)-t\,\nu(\Omega)>0$ となるように取り，
  補題~\ref{lem:rn-positive-set}を $\lambda:=\beta-t\nu$ に適用して $P$ を得る．
  $P$ の可測部分集合 $C$ について $t\,\nu(C)\leq\beta(C)$ であるから，
  任意の $B\in\mathcal{A}$ で
  \[
    \int_B(h+t\mathbf{1}_P)\,\d\nu
    =\int_Bh\,\d\nu+t\,\nu(B\cap P)
    \leq\int_Bh\,\d\nu+\beta(B\cap P)
    \leq\int_Bh\,\d\nu+\beta(B)=\alpha(B).
  \]
  よって $h+t\mathbf{1}_P\in\mathcal{G}$ である．
  一方，$\nu(P)=0$ なら仮定 2 から $\alpha(P)=0$，したがって $\beta(P)=0$ となり，
  $\lambda(P)=\beta(P)-t\nu(P)>0$ に反する．よって $\nu(P)>0$ であり，
  \[
    \int_\Omega(h+t\mathbf{1}_P)\,\d\nu=s+t\,\nu(P)>s
  \]
  となって $s$ の定義に反する．したがって $\beta(\Omega)=0$ であり，
  すべての $B$ で $\alpha(B)=\int_Bh\,\d\nu$ が成り立つ．
  一意性は補題~\ref{lem:rn-density-uniqueness}による．
\end{proof}

\begin{corollary}{密度による積分の書き換え}{rn-density-integral}
  定理~\ref{thm:meas-radon-nikodym}の条件 3 のもとで，
  任意の非負可測関数 $g\colon\Omega\to[0,\infty]$ に対し
  \[
    \int_\Omega g\,\d\alpha=\int_\Omega g\,h\,\d\nu
  \]
  が成り立つ．
\end{corollary}

\begin{proof}
  $g=\mathbf{1}_B$ のときは条件 3 そのものである．
  積分の線形性から非負単関数について成り立ち，
  非負単関数の列 $s_n\uparrow g$（注意~\ref{rem:meas-measurable-ops}）を取れば
  $s_nh\uparrow gh$ なので，両辺に単調収束定理（定理~\ref{thm:meas-mct}）を用いて一般の場合を得る．
\end{proof}
